\chapter{STRIDE Threat Model for Secure Online Banking Architecture}
% Running header: short title (TOC still uses full chapter title above)
\chaptermark{STRIDE Threat Model for Online Banking}
\label{ch:stride-banking}

This chapter presents a detailed STRIDE threat model for the secure online banking architecture. By decomposing the system into its core components---Frontend, Backend, and Data tiers---we identify potential vulnerabilities and map them to the security controls implemented in the project.

\section{System decomposition and trust boundaries}

The architecture is divided into distinct zones separated by \emph{trust boundaries}, where data transitions between different levels of control.

\subsection*{Public zone}
Includes Internet users (customers, tellers, admins) accessing the system via HTTPS/TLS~1.3.

\subsection*{DMZ (demilitarized zone)}
Houses the frontend Python Flask server and the Web Application Firewall (WAF).

\subsection*{Private application zone}
Contains the backend C\# .NET Core Web API, which handles business logic and transaction processing.

\subsection*{Restricted data zone}
The most secure tier, containing the Azure SQL Database and audit logs, protected by strict SQL firewalls.

\begin{figure}[htbp]
  \centering
  \includegraphics[width=\textwidth,height=0.88\textheight,keepaspectratio]{Blank_diagram.pdf}
  \caption{Trust boundaries and system decomposition (architecture overview).}
  \label{fig:architecture-zones}
\end{figure}

Figure~\ref{fig:architecture-zones} gives an overview of the zoning model and trust boundaries discussed above.

\section{STRIDE threat analysis and mitigations}

Table~\ref{tab:stride} analyzes the banking system's architecture through the STRIDE lens, identifying specific threats and the corresponding controls implemented to mitigate them.

\begin{table}[htbp]
  \centering
  \caption{STRIDE mapping: threats and implemented mitigations}
  \label{tab:stride}
  \small
  \setlength{\tabcolsep}{6pt}
  \renewcommand{\arraystretch}{1.25}
  \begin{tabularx}{\textwidth}{@{} >{\raggedright\arraybackslash}p{3.1cm} X X @{}}
    \toprule
    \textbf{Threat category} &
      \textbf{Potential project risk} &
      \textbf{Implemented \& compensating mitigations} \\
    \midrule
    Spoofing &
      Attackers stealing JWT tokens or bypassing MFA to impersonate a user. &
      JWT and TOTP MFA: JWT authentication and TOTP-based MFA; further mitigated by HttpOnly/Secure cookies and short token expiry. \\
    Tampering &
      Unauthorized modification of transaction data or interception of data in transit. &
      TLS~1.3 and hashing: All communication encrypted via TLS~1.3. User data protected using bcrypt/Argon2 with unique random salts. \\
    Repudiation &
      An admin or user performing an action and later denying it. &
      Audit logging: Separate audit servers with hash-protected integrity checks so logs are tamper-evident. \\
    Information disclosure &
      Exposure of sensitive banking data via man-in-the-middle (MITM) attacks. &
      Azure Key Vault: Centralized management of secrets and encryption keys. HSTS used to prevent TLS downgrade attacks. \\
    Denial of service &
      High-volume traffic or DDoS attacks overwhelming the banking API. &
      WAF and rate limiting: External firewalls provide DDoS mitigation; rate limiters at the edge restrict suspicious traffic spikes. \\
    Elevation of privilege &
      A customer attempting to access admin-level transaction controls. &
      RBAC: Strict role-based access control (customer, teller, admin) enforced by the backend API and monitored via audit logs. \\
    \bottomrule
  \end{tabularx}
\end{table}

\section{Defense-in-depth summary}

To address the identified threats, the project employs a defense-in-depth strategy:

\begin{description}
  \item[Network level:]
    Segregation via internal and external firewalls and a DMZ prevents direct access to the database from the Internet.
  \item[Application level:]
    The Web Application Firewall (WAF) protects against common exploits, complemented by robust input validation.
  \item[Identity level:]
    RBAC and managed identities ensure that even if one component is compromised, lateral movement is limited by least-privilege policies.
\end{description}

\section{Security gaps and continuous improvement}

Although the current architecture is robust, the threat model highlights areas for strengthening---for example, implementing HSTS to prevent protocol downgrades and ensuring secure MFA recovery processes for users who lose access to their devices.
